\clearpage
\noindent\makebox[10mm][l]{42}\parbox[t]{\dimexpr\textwidth-10mm\relax}{\centering\scriptsize BESTIMMUNG DES BREITENUNTERSCHIEDES ZWISCHEN GÖTTINGEN UND ALTONA.}\par
1820 von der untern Culmination des 13. Mai an gerechnet:
\begin{center}Mai 13.\quad Obere Culm.\ $-0''{,}10$;\end{center}
1824 von der obern Culmination des 20. April an gerechnet:
\[
\begin{array}{rl|r}
\text{Apr.}&21.\ \text{U. C.}&-0''{,}13\\
&21.\ \text{O. C.}&-0{,}26\\
&25.\ \text{U. C.}&-1{,}29\\
&27.\ \text{O. C.}&-2{,}04\\
&28.\ \text{O. C.}&-2{,}32\\
&29.\ \text{U. C.}&-2{,}45\\
\text{Mai}&1.\ \text{U. C.}&-2{,}93\\
&1.\ \text{O. C.}&-3{,}03\\
&2.\ \text{U. C.}&-3{,}14\\
&8.\ \text{O. C.}&-4{,}64\\
&9.\ \text{U. C.}&-4{,}77
\end{array}
\]
\begin{center}14.\end{center}
Bezeichnet man die Biegung des Fernrohrs, oder die Veränderung der Lage der auf die Ebene des getheilten Kreises projicirten optischen Axe gegen die Eintheilung, vermöge der Einwirkung der Schwere auf sämmtliche verbundene Bestandtheile des Instruments, bei horizontaler Lage der optischen Axe durch $f$, bei verticaler durch $g$, und setzt voraus, dass diese Biegung der Schwerkraft proportional ist (was bei der äusserst geringen Grösse der ganzen Wirkung unbedenklich scheint), so wird bei der Neigung der optischen Axe $z$ die Biegung durch $f\sin z+g\cos z$ ausgedrückt werden, so verstanden, dass wenn der Collimationsfehler $=e$ und die abgelesene Zenithdistanz $=z$ ist, die wahre Zenithdistanz
\[
=z-e+f\sin(z-e)+g\cos(z-e)
\]
sein wird. Wäre das Fernrohr vollkommen symmetrisch, so würde $g$ ganz wegfallen; allein da keine menschliche materielle Arbeit absolut vollkommen ist, und überdies die vollkommene Symmetrie schon durch die Balancirgewichte gewissermaassen gestört wird, so scheint es durchaus nicht ungereimt, die Möglichkeit eines ein oder ein Paar Zehntheile einer Secunde betragenden

\clearpage
\noindent\mbox{}\hfill III.\quad RESULTATE.\hfill 43\par
\noindent Werthes von $g$ zuzugeben, und wenn einmal die Rechnung auf einzelne Zehntheile oder gar Hunderttheile der Secunde genau geführt wird, so würde es inconsequent sein, die Berücksichtigung des zweiten Theils der Biegung, insofern sie möglich ist, zu unterlassen.
\begin{center}15.\end{center}
Das Complement des halben Unterschiedes der direct und durch Reflexion gemessenen Zenithdistanz zu $90^\circ$ gibt die Zenithdistanz vom Collimationsfehler und von dem ersten Theile der Biegung befreit, also bloss noch den zweiten Theil der Biegung enthaltend, und zwar mit entgegengesetztem Zeichen, [je] nachdem der Kreis in Osten oder Westen ist. Offenbar bezieht sich diese Zenithdistanz auf die Verticale an der Stelle, wo die optische Axe das Wassergefäss trifft, welche, für beide Culminationen des Nordsterns unmerklich verschieden, um $0''{,}05$ nördlicher ist als die Axe des Kreises. Diese Combination ist unserm Zweck auch insofern angemessener, als man der Voraussetzung der Unveränderlichkeit des Collimationsfehlers während der ganzen Dauer der Beobachtungen von 1824 ausweicht. Das Gewicht jener Bestimmung wird, wenn man die Anzahl der directen Beobachtungen $=\alpha$, die der Reflexionsbeobachtungen $=\beta$ setzt, $=\dfrac{4\alpha\beta}{\alpha+\beta}$, insofern man die Beobachtungsfehler als rein zufällig und von einander unabhängig betrachtet.
\begin{center}16.\end{center}
Bezeichnen wir nun mit
\begin{quote}
\begin{tabular}{@{}l@{\quad}p{.84\textwidth}@{}}
$\varphi$&die Polhöhe an dem Platz des Wassergefässes\\
$\delta$&die Declination des Nordsterns in der untern Culmination des 13. Mai 1820\\
$\delta'$&die Declination in der obern Culmination des 20. April 1824,
\end{tabular}
\end{quote}
so geben uns die Beobachtungen folgende Bestimmungen:
\begin{center}
\begin{tabular}{lllcll}
\multicolumn{6}{l}{für $\delta+\varphi-0{,}765g$}\\
\quad1820&Mai 13&$\ldots\,\ldotp$&$139^\circ52'38''{,}40$&Gewicht&$6{,}86$\\
\multicolumn{6}{l}{für $\delta-\varphi+0{,}800g$}\\
\quad1820&Mai 13&$\ldots\,\ldotp$&$36^\circ49'\phantom0 1''{,}50$&Gewicht&$2{,}00$
\end{tabular}
\end{center}
\noindent\hfill $6^*$\par

\clearpage
\noindent\makebox[10mm][l]{44}\parbox[t]{\dimexpr\textwidth-10mm\relax}{\centering\scriptsize BESTIMMUNG DES BREITENUNTERSCHIEDES ZWISCHEN GÖTTINGEN UND ALTONA.}\par
\begin{center}
\begin{tabular}{rrrc rr r c r}
\multicolumn{9}{l}{für $\delta'+\varphi-0{,}765g$}\\
1824&Apr.&21&$\ldots\,\ldotp$&$139^\circ$&$54'$&$5''{,}61$&Gewicht&$6{,}86$\\
&&25&$\ldots\,\ldotp$&139&54&$5{,}76$&»&$6{,}86$\\
&&29&$\ldots\,\ldotp$&139&54&$6{,}36$&»&$6{,}86$\\
&Mai&1&$\ldots\,\ldotp$&139&54&$5{,}91$&»&$6{,}86$\\
\multicolumn{9}{l}{für $\delta'-\varphi+0{,}800g$}\\
1824&Apr.&20&$\ldots\,\ldotp$&$36^\circ$&$50'$&$31''{,}15$&Gewicht&$2{,}67$\\
&&21&$\ldots\,\ldotp$&36&50&$30{,}29$&»&$6{,}86$\\
&&27&$\ldots\,\ldotp$&36&50&$30{,}65$&»&$6{,}86$\\
&&28&$\ldots\,\ldotp$&36&50&$30{,}73$&»&$6{,}86$\\
&Mai&1&$\ldots\,\ldotp$&36&50&$30{,}25$&»&$6{,}86$\\
\multicolumn{9}{l}{für $\delta'+\varphi+0{,}765g$}\\
1824&Mai&2&$\ldots\,\ldotp$&$139^\circ$&$54'$&$6''{,}71$&Gewicht&$6{,}86$\\
&&9&$\ldots\,\ldotp$&139&54&$6{,}15$&»&$6{,}86$\\
\multicolumn{9}{l}{für $\delta'-\varphi-0{,}800g$}\\
1824&Mai&8&$\ldots\,\ldotp$&$36^\circ$&$50'$&$30''{,}48$&Gewicht&$6{,}86$.
\end{tabular}
\end{center}
Wir erhalten demnach zur Bestimmung der vier unbekannten Grössen $\delta$, $\delta'$, $\varphi$, $g$ die sechs Gleichungen
\[
\begin{array}{rclrrrcr}
\delta+\varphi-0{,}765g&=&139^\circ&52'&38''{,}40&&\text{Gewicht}&6{,}86\\
\delta-\varphi+0{,}800g&=&36&49&1{,}50&&\text{»}&2{,}00\\
\delta'+\varphi-0{,}765g&=&139&54&5{,}91&&\text{»}&27{,}43\\
\delta'-\varphi+0{,}800g&=&36&50&30{,}54&&\text{»}&30{,}10\\
\delta'+\varphi+0{,}765g&=&139&54&6{,}43&&\text{»}&13{,}71\\
\delta'-\varphi-0{,}800g&=&36&50&30{,}48&&\text{»}&6{,}86.
\end{array}
\]
woraus sich durch die Methode der kleinsten Quadrate\textsuperscript{*)} folgende Werthe ergeben:
\[
\begin{array}{rcrrr}
\delta&=&88^\circ&20'&50''{,}33\\
\delta'&=&88&22&18{,}28\\
\varphi&=&51&31&47{,}90\\
g&=&\multicolumn{3}{r}{+\ 0{,}17.}
\end{array}
\]
Das Gewicht der Bestimmung von $\varphi$ wird hiebei $=60{,}8$.
\par\noindent\rule{.25\textwidth}{.4pt}\par
{\footnotesize *) Hier etwas bequemer nach dem Verfahren im Supplem.\ Theor.\ Comb.\ Observ.\par}

\clearpage
\noindent\mbox{}\hfill III.\quad RESULTATE.\hfill 45\par
Um für die Genauigkeit der Beobachtungen einigermaassen einen Maassstab zu haben, substituiren wir diese Werthe in den vierzehn Gleichungen, aus welchen die vorigen sechs zusammengezogen waren; es bleiben dann folgende Fehler übrig:
\begin{center}
\begin{tabular}{r|c}
Fehler&Gewicht der Gleichung\\\hline
$-0''{,}30$&$6{,}86$\\
$+1{,}07$&$2{,}00$\\
$+0{,}44$&$6{,}86$\\
$+0{,}29$&$6{,}86$\\
$-0{,}31$&$6{,}86$\\
$+0{,}14$&$6{,}86$\\
$-0{,}63$&$2{,}67$\\
$+0{,}23$&$6{,}86$\\
$-0{,}13$&$6{,}86$\\
$-0{,}21$&$6{,}86$\\
$+0{,}27$&$6{,}86$\\
$-0{,}40$&$6{,}86$\\
$+0{,}16$&$6{,}86$\\
$-0{,}24$&$6{,}86$
\end{tabular}
\end{center}
Die Summe der Producte der Quadrate dieser Fehler in die Gewichte wird $=9{,}6154$; also ein genäherter Werth für den mittlern Fehler einer Beobachtung
\[
=\sqrt{\frac{9{,}6154}{10}}=0''{,}981.
\]
Der mittlere in dem Endresultat für die Polhöhe zu befürchtende Fehler, so weit er von unregelmässig wirkenden Ursachen herrührt, ist demnach
\[
=\frac{0''{,}981}{\sqrt{60{,}8}}=0''{,}126.
\]
Etwas muss aber die Unsicherheit des Resultats allerdings grösser sein, da die Voraussetzung, dass sämmtliche Beobachtungsfehler ohne Ordnung von einander unabhängig sind, nicht ganz richtig ist. Bei gleichnamigen Beobach-

\clearpage
\noindent\makebox[10mm][l]{46}\parbox[t]{\dimexpr\textwidth-10mm\relax}{\centering\scriptsize BESTIMMUNG DES BREITENUNTERSCHIEDES ZWISCHEN GÖTTINGEN UND ALTONA.}\par
\noindent tungen zu einer Culmination, und bei den gleichnamigen Culminationen an mehrern Tagen liegt nemlich nahe dasselbe Ablesungsresultat zum Grunde, und obgleich, bei der Ablesung durch Verniers, fast immer andere Theilstriche sprechend werden, deren unregelmässige Theilungsfehler also bei unserm Verfahren in den mittlern Fehler einer Beobachtung $0''{,}981$ mit eingeschlossen sind, so ist doch natürlich, dass in den verschiedenen Gegenden des Limbus gewisse ungleiche Durchschnittsfehler vorherrschen müssen. Jedenfalls sind aber dieselben sehr klein. Im Jahre 1826 habe ich mit vier vortrefflichen Mikroskopen von \textsc{Repsold} 30 Theilstriche von 12 zu 12 Grad mit äusserster Sorgfalt geprüft, wobei jeder Theilstrich fast 200 mal, in abgeänderten Combinationen, eingestellt wurde. Das Resultat ist, dass das Mittel der Fehler von zwei diametral entgegengesetzten Theilstrichen, $A$ und $A+180^\circ$, so weit noch einige Regelmässigkeit zu erkennen ist, durch die Formel
\[
-1''{,}23\cos(2A-28^\circ28')-0''{,}22\cos(4A-47^\circ56')
\]
möglichst nahe dargestellt wird, dass die dann übrig bleibenden Fehler als regellos erscheinen, und die Quadratwurzel aus dem Mittel ihrer Quadrate $=0''{,}32$ wird. Ich hatte mir vorgesetzt, diese Prüfung auf die doppelte Anzahl der Theilstriche auszudehnen; allein bei der Geringfügigkeit der sich ergebenden Resultate scheint diese Untersuchung den grossen dazu erforderlichen Zeitaufwand nicht zu verdienen. Es bedarf keiner Erinnerung, dass der erste Theil des regelmässigen Fehlers $-1''{,}23\cos(2A-28^\circ28')$ von selbst wegfällt, wenn, wie bei obigen Beobachtungen immer geschehen ist, alle vier Verniers abgelesen werden. Er enthält hingegen eine reelle Verbesserung, falls man die Theilung nur an zwei gegenüberliegenden Stellen abliest, wie ich gegenwärtig immer thue, seitdem ich mich mit bedeutendem Gewinn für die Feinheit der Ablesung statt der Verniers zweier \textsc{Repsold}scher Mikroskope bediene.
\begin{center}17.\end{center}
Zieht man vor, $g=0$ vorauszusetzen, so fällt die Polhöhe um $0''{,}07$ kleiner aus, und das Gewicht dieser Bestimmung wird $=84{,}1$. Anderweitige, an einem andern Orte anzuführende Beobachtungen scheinen übrigens den obigen

\clearpage
\noindent\mbox{}\hfill III.\quad RESULTATE.\hfill 47\par
\noindent Werth von $g$, dem Zeichen und auch sehr nahe der Grösse nach, zu bestätigen, reichen aber noch nicht hin, über einen so delicaten Gegenstand zu entscheiden.

Den Coefficienten $f$ kann man aus vorliegenden Beobachtungen nicht bestimmen, ohne die Unveränderlichkeit des Collimationsfehlers während der Beobachtungen von 1824 vorauszusetzen. Erlaubt man sich diese Voraussetzung, so hat man 28 Gleichungen, deren gehörige Behandlung
\[
\begin{array}{rcll}
\varphi&=&51^\circ31'47''{,}89&\text{mit dem Gewicht }60{,}9\\
f&=&\phantom{51^\circ31'}+\ 0{,}76&\\
g&=&\phantom{51^\circ31'}+\ 0{,}23&
\end{array}
\]
gibt. Da man gegenwärtig, durch Einstellen des Fernrohrs auf den Nadirpunkt, den Collimationsfehler jede Stunde mit bewundernswürdiger Genauigkeit ohne Umlegen bestimmen kann\textsuperscript{*)}, so behalte ich mir weitere Prüfung dieses Gegenstandes vor.
\begin{center}18.\end{center}
Mit Vorbehalt der durch künftige weitere Untersuchungen noch auszumittelnden Correction, die wohl schwerlich eine halbe Secunde erreichen kann, setze ich daher die Polhöhe
\begin{center}
\begin{tabular}{p{.72\textwidth}r}
in Göttingen&\\
\quad für den Platz des Wassergefässes bei den Nordsternbeobachtungen&$51^\circ31'47''{,}90$\\
\quad für den Platz des \textsc{Reichenbach}schen Meridiankreises&$47{,}85$\\
\quad für den Platz des Zenithsectors&$47{,}92$\\
\multicolumn{2}{p{.94\textwidth}}{\quad\quad (welche letztere zur Reduction der Declinationen der Zenithalsterne zum Grunde gelegt ist)}\\
in Altona&\\
\quad für den Platz des Zenithsectors&$53^\circ32'44''{,}42$\\
\quad für den Platz des Meridiankreises&$45{,}27$.
\end{tabular}
\end{center}
\par\noindent\rule{.25\textwidth}{.4pt}\par
{\footnotesize *) Ich bediene mich dieses unschätzbaren Mittels, dessen Ausführbarkeit \textsc{Bohnenberger} zuerst gezeigt hat, seit zwei Jahren beständig.\par}

\clearpage
\noindent\makebox[10mm][l]{48}\parbox[t]{\dimexpr\textwidth-10mm\relax}{\centering\scriptsize BESTIMMUNG DES BREITENUNTERSCHIEDES ZWISCHEN GÖTTINGEN UND ALTONA.}\par
\begin{center}19.\end{center}
Nach der trigonometrischen Verbindung der Sternwarten von Göttingen und Altona liegt letztere
\[
\begin{array}{rl}
115163{,}725&\text{Toisen nördlich}\\
7{,}211&\text{Toisen westlich}
\end{array}
\]
von jener. Diese Zahlen beziehen sich auf die Plätze der Meridiankreise; sie gründen sich auf den Werth der Dreiecksseite Hamburg-Hohenhorn $13841{,}815$ Toisen, und diese auf die von Hrn.\ Prof.\ \textsc{Schumacher} in Holstein im Jahr 1820 gemessene Basis. Da jedoch die Vergleichung der dabei gebrauchten Messstangen mit der Normaltoise noch nicht \textls[150]{definitiv} vollendet ist, so wird obige Entfernung in Zukunft noch in demselben Verhältniss abzuändern sein, wie die Basis selbst, welche Veränderung aber jedenfalls nur sehr gering sein kann. Der mittlere Breitengrad zwischen beiden Sternwarten ergibt sich danach
\[
=57127{,}2\text{ Toisen},
\]
merklich grösser als man nach den mittlern Werthen der in Frankreich und England gemessenen Grade hätte erwarten sollen.
\begin{center}20.\end{center}
Die hannoversche Gradmessung liefert also einen neuen Beitrag zur Bestätigung der nicht mehr zu bezweifelnden Wahrheit, dass die Oberfläche der Erde keine ganz regelmässige Gestalt hat. Von dieser Unregelmässigkeit haben bereits die Anomalien bei den Theilen der französischen und der englischen Gradmessung Beweise gegeben, noch stärkere die Anomalien bei den Polhöhen mehrerer Örter in Italien. Bei der hannoverschen Gradmessung findet sich ausser der Anomalie zwischen Göttingen und Altona eine noch beträchtlich stärkere bei einem zwischenliegenden Dreieckspunkte, dem Brocken. Wenn man meine Dreiecke als auf der Oberfläche eines elliptischen Sphäroids liegend, dessen Dimensionen die von \textsc{Walbeck} aus der Gesammtheit der bisherigen Gradmessungen abgeleiteten sind, und welches nach unserer besten gegenwärtigen Kenntniss sich am vollkommensten an die wirkliche Gestalt

\clearpage
\noindent\mbox{}\hfill III.\quad RESULTATE.\hfill 49\par
\noindent\textls[150]{im Ganzen} anschliesst (Abplattung $\dfrac1{302{,}78}$, der dreihundertsechzigste Theil des Erdmeridians $=57009{,}758$ Toisen), berechnet, und dabei von der Polhöhe von Göttingen $=51^\circ31'47''{,}85$ ausgeht, so findet sich die Breite
\[
\begin{array}{lcrrr}
\text{des Brockens}&=&51^\circ&48'&1''{,}85\\
\text{von Altona}&=&53&32&50{,}79.
\end{array}
\]
Während nun die astronomischen Beobachtungen die Polhöhe von Altona $5''{,}52$ \textls[150]{kleiner} gegeben haben, geben die von Hrn.\ von \textsc{Zach} auf dem Brocken angestellten Beobachtungen die Polhöhe dieses Punktes $10\text{—}11''$ \textls[150]{grösser}\textsuperscript{*)}, ein Unterschied, von dem doch jedenfalls nur ein kleiner Theil dem Instrumente und den in der Rechnung gebrauchten Declinationen zur Last fallen kann. Die Vergleichung des Breitenunterschiedes zwischen Altona und dem Brocken mit der Krümmung, welche dem sich der Erde im Ganzen am besten anschliessenden Sphäroid entspricht, würde daher eine Abweichung von $16''$ geben.

Nach unserm Dafürhalten betrachtet man diesen Gegenstand aus einem falschen Gesichtspunkte, wenn man bei solchen Erscheinungen immer nur von Localablenkungen der Lothlinie spricht, und sie also gleichsam nur als einzelne Ausnahmen ansieht. Was wir im geometrischen Sinn Oberfläche der Erde nennen, ist nichts anderes als diejenige Fläche, welche überall die Richtung der Schwere senkrecht schneidet, und von der die Oberfläche des Weltmeers einen Theil ausmacht. Die Richtung der Schwere an jedem Punkte wird aber durch die Gestalt des festen Theils der Erde und seine ungleiche Dichtigkeit bestimmt, und an der äussern Rinde der Erde, von der allein wir etwas wissen, zeigt sich diese Gestalt und Dichtigkeit als höchst unregelmässig; die Unregelmässigkeit der Dichtigkeit mag sich leicht noch ziemlich tief unter die äussere Rinde erstrecken, und entzieht sich ganz unsern Berechnungen, zu welchen fast alle Data fehlen. Die geometrische Oberfläche ist das Product der Gesammtwirkung dieser ungleich vertheilten Elemente, und anstatt vorkommende unzweideutige Beweise der Unregelmässigkeit befremdend zu finden, scheint es eher zu bewundern, dass sie nicht noch grösser ist.
\par\noindent\rule{.25\textwidth}{.4pt}\par
{\footnotesize *) Monatl.\ Corresp.\ B.\ X.\ S.\ 203. An einem Platze, der etwa $0''{,}5$ südlicher liegt, als der Dreieckspunkt, fand dieser geschickte Beobachter aus 188 Beobachtungen von $\alpha$ Aquilae $51^\circ48'12''{,}12$. Aus Sonnenbeobachtungen fand er $51^\circ48'11''{,}17$.\par}
\noindent IX.\hfill 7\par

\clearpage
\noindent\makebox[10mm][l]{50}\parbox[t]{\dimexpr\textwidth-10mm\relax}{\centering\scriptsize BESTIMMUNG DES BREITENUNTERSCHIEDES ZWISCHEN GÖTTINGEN UND ALTONA.}\par
\noindent Wären die astronomischen Beobachtungen einer zehn- oder hundertmal grössern Genauigkeit fähig, als sie gegenwärtig haben, so würden sie diese Unregelmässigkeit ohne Zweifel überall nachweisen.

Bei dieser Lage der Sache hindert aber noch nichts, die Erde im Ganzen als ein elliptisches Revolutionssphäroid zu betrachten, von dem die wirkliche (geometrische) Oberfläche überall bald in stärkern, bald in schwächern, bald in kürzern, bald in längern Undulationen abweicht. Wäre es möglich, die ganze Erde mit Einem trigonometrischen Netze gleichsam zu umspinnen, und die gegenseitige Lage aller Punkte dadurch zu berechnen, so würde das idealische Revolutionssphäroid dasjenige sein, auf welchem berechnet die Richtungen der Verticalen die möglich beste Übereinstimmung mit den astronomischen Beobachtungen gäben. Wenn man gleich von diesem unerreichbaren Ideale immer weit entfernt bleiben wird, so leidet es doch keinen Zweifel, dass die künftigen Jahrhunderte die mathematische Kenntniss der Erdfigur sehr viel werden weiter bringen können. Die Vervielfältigung der Gradmessungen ist aber eigentlich nur der Anfang dazu, woraus nur einzelne Resultate für eine kleine Anzahl in isolirten Linien liegender Punkte hervorgehen: wie viel ergiebiger wird aber die Ausbeute sein, wenn diejenigen trigonometrischen Operationen, welche mit ausgesuchten Hülfsmitteln in verschiedenen Ländern ausgeführt sind, in Verknüpfung kommen und sich zu Einem grossen System abrunden. Vielleicht ist die Aussicht nicht chimärisch, dass einst alle Sternwarten von Europa trigonometrisch unter einander verbunden sein werden, da schon jetzt solche Verbindungen von Schottland bis zum adriatischen Meere und von Formentera bis Fünen vorhanden, wenn gleich bisher nur theilweise öffentlich bekannt gemacht sind. Möchte nur dieser letzte Umstand, mehr als bisher geschehen, beachtet, und kostbare Materialien, die der wissenschaftlichen Welt angehören sollten, dieser nicht entzogen, oder gar der Gefahr des Untergangs preisgegeben werden!
\begin{center}21.\end{center}
Ein nicht uninteressantes Resultat gibt noch die Vergleichung der aus den Sectorbeobachtungen hervorgegangenen Sterndeclinationen mit ältern Bestimmungen, wo solche vorhanden sind. Von unsern 43 Sternen finden sich

\clearpage
\noindent\mbox{}\hfill III.\quad RESULTATE.\hfill 51\par
\noindent 27 in \textsc{Piazzi}s und 13 in \textsc{Bessel}s \textsc{Bradley}schem Catalog. Hier folgt die Vergleichung unserer Bestimmungen (1827) mit \raisebox{.3ex}{.}den \textsc{Bradley}schen (1755) und \textsc{Piazzi}schen (1800), nach \textsc{Bessel}s neuer Bestimmung der Präcession reducirt; positive Zeichen bedeuten eine nördlichere Stellung aus unserer Bestimmung.
\begin{center}
{\small\setlength{\tabcolsep}{4pt}
\begin{tabular}{r|l|r|r||r|l|r|r}
\multicolumn{2}{c|}{Bezeichnung}&\textsc{Bradley}&\textsc{Piazzi}&\multicolumn{2}{c|}{Bezeichnung}&\textsc{Bradley}&\textsc{Piazzi}\\\cline{1-4}\cline{5-8}
1&24 Canum&$+0''{,}2$&$+1''{,}1$&17&P.\ 15.39&—&$-0''{,}7$\\
2&83 Ursae&$-1{,}5$&$-2{,}0$&25&$\theta$ Draconis&$+23''{,}5$&$+8{,}6$\\
3&$\eta$ Ursae&$-2{,}4$&$-2{,}3$&27&P.\ 16.33&—&$-2{,}6$\\
4&86 Ursae&$-5{,}4$&$-0{,}8$&28&P.\ 16.56&—&$-2{,}6$\\
6&P.\ 13.289&—&$-1{,}3$&32&16 Draconis&$+1{,}2$&$-3{,}7$\\
7&13 Bootis&$+1{,}4$&$+2{,}4$&36&P.\ 16.253&—&$-2{,}1$\\
8&$\varkappa$ Bootis sq.&$-2{,}9$&$-2{,}2$&37&P.\ 16.291&—&$+10{,}3$\\
9&P.\ 14.56&—&$-0{,}4$&38&P.\ 16.310&—&$-4{,}9$\\
10&$\theta$ Bootis&$-30{,}6$&$-10{,}8$&39&P.\ 17.20&—&$-3{,}0$\\
11&P.\ 14.131&—&$+7{,}4$&40&P.\ 17.38&—&$-2{,}9$\\
12&P.\ 14.164&—&$+0{,}1$&41&74 Herculis&$+1{,}7$&$+0{,}7$\\
13&39 Bootis med.&$+4{,}8$&$+2{,}7$&42&P.\ 17.120&—&$-2{,}1$\\
14&P.\ 14.235&—&$-5{,}0$&43&$\beta$ Draconis&$-0{,}5$&$-1{,}5$\\
15&44 Bootis med.&$+1{,}9$&$+1{,}1$&&&&
\end{tabular}}
\end{center}
\vspace{14mm}
\begin{center}\rule{.25\textwidth}{.4pt}\end{center}
\noindent\hfill $7^*$\par
